\subsection{Simultaneous spectrum}

Let \(\Lambda\) a set. Let \(C_{\Lambda} = \mathbf{C}[(X_{\lambda})_{\lambda \in \Lambda}]\) be the unital algebra of complex polynomials in a family of indeterminates \((X_{\lambda})_{\lambda \in \Lambda}\) . For every \(\chi \in X(C_{\Lambda})\) , one has \((\chi (X_{\lambda}))_{\lambda \in \Lambda} \in C^{\Lambda}\) ; the map \(\chi \mapsto (\chi (X_{\lambda}))_{\lambda \in \Lambda}\) is a homeomorphism from \(X(C_{\Lambda})\) onto the product space \(C^{\Lambda}\) , by which these spaces are identified.

On the other hand, let A be a commutative unital Banach algebra and let \(x = (x_{\lambda})_{\lambda \in \Lambda}\) be a family of elements of A. There exists a unique unital morphism h from \(C_{\Lambda}\) to A such that \(h(\mathrm{X}_{\lambda}) = x_{\lambda}\) for all \(\lambda\) . The continuous map \(\mathsf{X}(h)\) of \(\mathsf{X}(\mathrm{A})\) into \(C^{\Lambda}\) is the map which to \(\chi\) associates the family \((\chi(x_{\lambda}))_{\lambda \in \Lambda}\) . It is called the map of \(\mathsf{X}(\mathrm{A})\) into \(C^{\Lambda}\) defined by x.

DEFINITION 2. --- The image of the map X(h) is called the simultaneous spectrum of x, and is denoted \(\mathrm{Sp}_{\mathrm{A}}^{\Lambda}(x)\) or \(\mathrm{Sp}^{\Lambda}(x)\) .

The simultaneous spectrum of \(x\) is a compact subset of \(\mathbf{C}^{\Lambda}\) . An element \(c = (c_{\lambda}) \in \mathbf{C}^{\Lambda}\) belongs to \(\mathrm{Sp}_{\mathrm{A}}^{\Lambda}(x)\) if and only if the elements \(x_{\lambda} - c_{\lambda}\) belong to a same maximal ideal of \(\mathrm{A}\) , in other words if the family \((x_{\lambda} - c_{\lambda})_{\lambda \in \Lambda}\) does not generate the algebra \(\mathrm{A}\) .

If \(\Lambda\) contains a single element, so that the family \(x\) reduces to a single element \(x \in \mathsf{A}\) , one has \(\mathrm{Sp}_{\mathsf{A}}^{\Lambda}(x) = \mathrm{Sp}_{\mathsf{A}}(x)\) (I, p.~37, prop. 6, \(b\) )). If \(\Lambda' \subset \Lambda\) , then \(\mathrm{Sp}_{\mathsf{A}}^{\Lambda'}((x_{\lambda})_{\lambda \in \Lambda'})\) is the image of \(\mathrm{Sp}_{\mathsf{A}}^{\Lambda}((x_{\lambda})_{\lambda \in \Lambda})\) by the canonical projection map of \(\mathbf{C}^{\Lambda}\) on \(\mathbf{C}^{\Lambda'}\) . In particular, we have

\[
\operatorname{Sp} _ {\mathrm{A}} ^ {\Lambda} (x) \subset \prod_ {\lambda \in \Lambda} \operatorname{Sp} _ {\mathrm{A}} (x _ {\lambda}).
\]

We denote \(z_{\lambda}\) , for \(\lambda \in \Lambda\) , the coordinate functions on \(\mathbf{C}^{\Lambda}\) . If \(\chi \in \mathsf{X}(\mathsf{A})\) , the value at \(\chi\) of \(z_{\lambda} \circ \mathsf{X}(h)\) is \(\chi(x_{\lambda})\) , hence \(z_{\lambda} \circ \mathsf{X}(h) = \mathcal{G}(x_{\lambda})\) .

Let A and B be commutative unital Banach algebras, \(\varphi\) a unital morphism from A to B, and \(x = (x_{\lambda})_{\lambda \in \Lambda}\) a family of elements of A. Let \(\varphi(x)\) the family \((\varphi(x_{\lambda}))_{\lambda \in \Lambda}\) of elements of B. We have, for every \(\chi \in \mathsf{X}(\mathsf{B})\) , and every \(\lambda \in \Lambda\)

\[
\chi (\varphi (x _ {\lambda})) = (\mathsf {X} (\varphi) (\chi)) (x _ {\lambda}),
\]

hence \(\operatorname{Sp}_{\mathrm{B}}^{\Lambda}(\varphi(x)) \subset \operatorname{Sp}_{\mathrm{A}}^{\Lambda}(x)\) . The diagram

\[
\begin{array}{c} \mathsf {X} (\mathrm{B}) \xrightarrow {\mathsf {X} (\varphi)} \mathsf {X} (\mathrm{A}) \\ \Big \downarrow \qquad \qquad \qquad \qquad \Big \downarrow \\ \operatorname{Sp} _ {\mathrm{B}} ^ {\Lambda} (\varphi (x)) \xrightarrow {i} \operatorname{Sp} _ {\mathrm{A}} ^ {\Lambda} (x) \end{array}\tag{1}
\]

where i denotes the inclusion, and the vertical arrows denote the maps defined by the families \(\varphi(x)\) and x, is therefore commutative.

Example. --- Let \(K \subset C^{\Lambda}\) a compact subset. Let \(z = (z_{\lambda})_{\lambda \in \Lambda}\) the family in \(\mathcal{C}(K)\) of the restrictions to K of the coordinate functions of \(\mathbf{C}^{\Lambda}\) . Then the simultaneous spectrum \(\mathrm{Sp}_{\mathcal{C}(\mathrm{K})}^{\Lambda}(z)\) is equal to K. Indeed, by cor. 2 of I, p.~32, every character \(\chi\) of \(\mathcal{C}(\mathrm{K})\) is of the form \(f \mapsto f(x)\) for an element \(x \in \mathrm{K}\) , and one then has \((\chi(z_{\lambda}))_{\lambda \in \Lambda} = x\) .

PROPOSITION 12. — Let \(\Lambda\) a set and A a commutative unital Banach algebra. Let \(x = (x_{\lambda})_{\lambda \in \Lambda}\) be a family of elements of A.

\begin{enumerate}
\def\labelenumi{\alph{enumi})}
\item
  We assume that the full subalgebra of A generated by the family x is dense in A. The map from X(A) to \(\mathbf{C}^{\Lambda}\) defined by x is a homeomorphism from X(A) onto the simultaneous spectrum \(\mathrm{Sp}_{\mathrm{A}}^{\Lambda}(x)\) ;
\item
  We assume that the unital subalgebra of A generated by the family x is dense in A. For every \(c \in C^{\Lambda}\) , the following conditions are equivalent:
\end{enumerate}

\begin{enumerate}
\def\labelenumi{(\roman{enumi})}
\item
  \(c \in \mathrm{Sp}_{\mathrm{A}}^{\Lambda}(x)\) ;
\item
  \(|\mathrm{P}(c)| \leqslant \varrho(\mathrm{P}(x))\) for every polynomial \(\mathrm{P} \in \mathbf{C}[(\mathrm{X}_{\lambda})_{\lambda \in \Lambda}]\) ;
\item
  \(|\mathrm{P}(c)| \leqslant \| \mathrm{P}(x)\|\) for every polynomial \(\mathrm{P} \in \mathbf{C}[(\mathrm{X}_{\lambda})_{\lambda \in \Lambda}]\) .
\end{enumerate}

\begin{enumerate}
\def\labelenumi{\alph{enumi})}
\item
  The map from \(\mathsf{X}(\mathsf{A})\) into \(\mathrm{Sp}_{\mathrm{A}}^{\Lambda}(x)\) defined by the family \(x\) is continuous and surjective. Let \(\chi\) , \(\chi' \in \mathsf{X}(\mathsf{A})\) be characters having the same image, that is, such that \(\chi(x_{\lambda}) = \chi'(x_{\lambda})\) for all \(\lambda \in \Lambda\) . The characters \(\chi\) and \(\chi'\) coincide on elements of the form \(\mathrm{P}(x)\mathrm{Q}(x)^{-1}\) , where \(\mathrm{P} \in \mathbf{C}[(\mathrm{X}_{\lambda})]\) , \(\mathrm{Q} \in \mathbf{C}[(\mathrm{X}_{\lambda})]\) and \(\mathrm{Q}(x)\) is invertible in A, that is, on the full subalgebra of A generated by the elements \(x_{\lambda}\) (lemma 2 of I, p.~6). Since \(\chi\) and \(\chi'\) are continuous (th. 1 of I, p.~29), they are therefore equal on A. This proves that \(\mathsf{X}(h)\) is a continuous bijection from \(\mathsf{X}(\mathsf{A})\) on \(\mathrm{Sp}_{\mathrm{A}}^{\Lambda}(x)\) , and consequently it is a homeomorphism since \(\mathsf{X}(\mathsf{A})\) is compact.
\item
  Let us show that (i) implies (ii): if \(c = (c_{\lambda})_{\lambda \in \Lambda} \in \mathrm{Sp}_{\mathrm{A}}^{\Lambda}(x)\) , there exists \(\chi \in \mathsf{X}(\mathsf{A})\) such that \(c_{\lambda} = \chi(x_{\lambda})\) for all \(\lambda\) . For every \(\mathrm{P} \in \mathbf{C}[(\mathrm{X}_{\lambda})]\) , we therefore have \(|\mathrm{P}(c)| = |\mathrm{P}((\chi(x_{\lambda}))_{\lambda \in \Lambda})| = |\chi(\mathrm{P}(x))| \leqslant \varrho(\mathrm{P}(x))\) .
\end{enumerate}

Assertion (ii) implies (iii) by virtue of the inequality \(\varrho(x) \leqslant \|x\|\) , valid for all \(x \in A\) by definition of the spectral radius.

Let us finally show that (iii) implies (i). Let \(c = (c_{\lambda})_{\lambda \in \Lambda} \in \mathbf{C}^{\Lambda}\) such that

\[
| \mathrm{P} (c) | \leqslant \| \mathrm{P} (x) \|\tag{2}
\]

for all \(P \in \mathbf{C}[(X_{\lambda})]\) . Let A' be the unital subalgebra of A generated by the family x; its elements are of the form \(P(x)\) for \(P \in \mathbf{C}[(X_{\lambda})]\) . The estimate (2) implies that the condition \(P(x) = 0\) implies \(\mathrm{P}(c)=0\) There therefore exists a unital algebra morphism \(\xi\) of \(A'\) into C such that \(\xi(x_{\lambda})=c_{\lambda}\) for all \(\lambda\in\Lambda\) . By (2), the morphism \(\xi\) is continuous. It therefore extends by continuity to a character \(\chi\) of \(\overline{\mathrm{A}'}=\mathrm{A}\) , which satisfies \(c=(\chi(x_{\lambda}))_{\lambda\in\Lambda}\in\mathrm{Sp}_{\mathrm{A}}^{\Lambda}(x)\) . This completes the proof.

